% © 2026 Ing. David Jaroš — CC BY-NC-ND 4.0
% =====================================================================
% RESEARCH TRACK — G137-139-METRIC-BRIDGE
% Target: research_tracks/sector_coupling/G137_139_metric_bridge.tex
% Status: §3 [DERIVED — reduced model], §4 [DERIVED — reduced model],
%         §5 [OPEN], §6 [OPEN]
% Layer: [L1 cond.] — biquaternionic geometry, 2D scalar reduction;
%        conditional on lifts listed in §6.
% =====================================================================

\section{Metric-Induced Coupling Between Winding Sectors 137 and 139}
\label{sec:g137_139_bridge}

\noindent\textbf{Status summary:}
\begin{itemize}
  \item Direct overlap $\langle 139|137\rangle = 0$: \textbf{NO-GO} (Fourier orthogonality on $S^1_\psi$).
  \item Common-multiple mode $137\cdot139$: \textbf{NO-GO} (divisibility of the index is not overlap).
  \item Explicit canonical $\Theta^3/\Theta^4$ vertex: \textbf{not present} in canon
        (EOM $\nabla^\dagger\nabla\Theta = \kappa\mathcal{T}$ is linear in $\Theta$ at fixed metric).
  \item Metric-induced $\Delta n = \pm 1$ coupling: \textbf{DERIVED} in the reduced model of §2–§4,
        and shown \emph{diffeomorphism-invariant} (§4).
  \item Direct $\Delta n = \pm 2$ element: vanishes at $O(\varepsilon)$ (§4) — the $137\to139$
        transition is second order, via the virtual (non-stable) intermediate $n=138$.
  \item No-communication: \textbf{unaffected}. The coupling is a causal, local interaction
        (mode conversion), not nonlocal signalling. Nothing here revives entanglement-based
        communication; correlations remain non-signalling.
\end{itemize}

% ---------------------------------------------------------------
\subsection{Input from canon}
\label{sec:bridge_input}

Two canonical ingredients only:
\begin{enumerate}
  \item The two-mode winding vacuum of
        \texttt{canonical/geometry/biquaternionic\_vacuum\_solutions.tex}:
        \begin{equation}
          \Theta_0(\psi) = A_1\,e^{i\psi/R_\psi} + A_2\,e^{2i\psi/R_\psi},
          \qquad A_1, A_2 \in \mathbb{B},
          \label{eq:bridge_vacuum}
        \end{equation}
        with $h_{\psi\psi} \neq 0$ [L1] and required for the physicality of the
        $\varphi$-universe parameter.
  \item The complex-time signature relation
        (\texttt{canonical/gr\_closure/step3\_signature\_theorem.tex}):
        $g_{tt} = -|\partial_\psi\Theta_0|^2$, so the \emph{same} profile
        $f(\psi) := |\partial_\psi\Theta_0|^2$ enters both $g_{\psi\psi}$ and $g_{tt}$.
\end{enumerate}

\noindent\textbf{Reduced model} [declared reductions, lifted in §6]:
complex scalar $\theta(t,\psi)$ instead of full biquaternion; 2D $(t,\psi)$;
source-free; $R_\psi = 1$; $\varepsilon := |A_2|/|A_1|$ perturbative.

% ---------------------------------------------------------------
\subsection{Naive mixing (coordinate-dependent)}
\label{sec:bridge_naive}

With $f(\psi) = A_1^2 + 4A_1\varepsilon\,(a_{2r}\cos\psi - a_{2i}\sin\psi) + O(\varepsilon^2)$,
the Laplace--Beltrami operator on the circle gives
\begin{equation}
  (L)_{n+1,\,n} = \frac{\varepsilon\,n\,(2n+1)\,(a_{2r} + i\,a_{2i})}{A_1^3} \neq 0.
  \quad\textbf{[DERIVED — reduced model; coordinate-dependent]}
\end{equation}
\textbf{Warning:} in pure 1D this is removable by $\psi$-reparametrization
(flattening $g_{\psi\psi}$), hence by itself it is \emph{not} evidence of physical mixing.

% ---------------------------------------------------------------
\subsection{Diffeomorphism test: the lapse profile survives}
\label{sec:bridge_diffeo}

Flatten $g_{\psi\psi}$ by $d\psi' = \sqrt{f}\,d\psi / c$, $c = \langle\sqrt{f}\rangle = A_1$.
The residual metric is
\begin{equation}
  ds^2 = -f\bigl(\psi(\psi')\bigr)\,dt^2 + c^2\,d\psi'^2 ,
\end{equation}
and the lapse profile $f(\psi(\psi'))$ retains its first harmonic: one reparametrization
function cannot flatten both $g_{\psi\psi}$ and $g_{tt}$. For modes
$e^{-i\omega t + in\psi'}$ the wave operator gives
\begin{equation}
  \boxed{\;
  (L)_{n,\,n} = \frac{\omega^2 - n^2}{A_1^2},\qquad
  (L)_{n+1,\,n} = -\,i\,\varepsilon\,\frac{(a_{2i} - i\,a_{2r})\,(n + 2\omega^2)}{A_1^3},\qquad
  (L)_{n+2,\,n} = 0 \;}
  \label{eq:bridge_invariant}
\end{equation}
\textbf{[DERIVED — reduced model; diffeo-invariant gauge]}

\noindent Consequences:
\begin{itemize}
  \item The $\Delta n = \pm1$ coupling is \emph{physical} (survives metric flattening;
        carried by the $g_{tt}$ lapse harmonic mandated by the signature theorem).
  \item Coupling strength grows as $\omega^2$: high-frequency driving couples
        sectors more strongly.
  \item $\Delta n = \pm2$ vanishes at first order: no direct $137\to139$.
\end{itemize}

% ---------------------------------------------------------------
\subsection{Second-order $137 \to 139$ amplitude}
\label{sec:bridge_second_order}

Via the virtual intermediate $n = 138$ (which need \emph{not} be a stable sector:
the Fourier tower contains all $n \in \mathbb{Z}$; the stable set
$S = \{2,127,137,139,151,157\}$ labels minima of $V_{\mathrm{eff}}$, not the spectrum):
\begin{equation}
  \mathcal{A}^{(2)}(137\to139)
  = \frac{(L)_{139,138}\,(L)_{138,137}}{\Delta E_{138}}
  \;\xrightarrow{\;\omega^2 = 137^2\;}\;
  \frac{5\,161\,612\;\varepsilon^2\,|A_2|^2}{A_1^4}.
  \quad\textbf{[DERIVED — reduced model]}
\end{equation}
Nonzero for any $A_2 \neq 0$. Physical meaning: a \emph{causal pumping channel}
$137 \to 139$ exists in the two-mode vacuum; sector 139 is dynamically addressable
by driving the 137 sector. This is mode conversion, not signalling.

% ---------------------------------------------------------------
\subsection{Charge consistency lemma [OPEN — blocking]}
\label{sec:bridge_charge}

In canon the winding $n$ on $S^1_\psi$ encodes electric charge. Then:
(i) $\Delta n = \pm1$ conversion must be absorbed by the background condensate;
(ii) a vacuum with coherent charged modes $n=1,2$ is a charged condensate,
i.e., spontaneously broken U(1) --- a superconducting vacuum that would Higgs
the photon. Required resolution (one of):
\begin{itemize}
  \item $n_\psi \neq Q_{\mathrm{EM}}$ directly (pre-projection topological index only), or
  \item the physical vacuum carries no net charged condensate (constrains $A_1, A_2$).
\end{itemize}
\textbf{Until resolved, the bridge is [L1 cond.] at best.} Note this issue predates
the bridge: it already threatens the physicality argument for $\varphi$.

% ---------------------------------------------------------------
\subsection{Open lifts and one canon conflict}
\label{sec:bridge_open}

\begin{enumerate}
  \item Lift scalar $\to$ full biquaternion (Sc$(\cdot)$ algebra; check no cancellation).
  \item Lift 2D $\to$ 4D+$\psi$ (mixed components $g_{\mu\psi}$ add channels, cannot remove them).
  \item Charge consistency lemma (§5).
  \item \textbf{Canon conflict to resolve:} \texttt{Rpsi\_dynamical\_fix.tex} derives
        $R_\psi = R_t$ at the self-dual point; \texttt{biquaternionic\_vacuum\_solutions.tex}
        \S{}Candidate~1 marks self-duality \textbf{DEAD END} ($\alpha = 1$ unphysical) and
        records $R_\psi = \hbar/(m_e c)$ as [KALIBROVÁNO]. Either two different meanings
        of $R_t$ (spatial torus vs.\ time circle) --- then both files must say so explicitly ---
        or a genuine contradiction.
\end{enumerate}

\noindent\textbf{Scope note.} This track establishes (conditionally) a causal inter-sector
coupling. It does not establish, and cannot establish, nonlocal communication,
readout of source-free configurations, or any claim quarantined under
\texttt{speculative\_extensions/}. See \texttt{docs/CONSCIOUSNESS\_CLAIMS\_ETHICS.md}.
